\subsection{Trace and Determinant}

\begin{proposition} \label{prop:trace-arithmetic} \hlt{Trace Arithmetic}\\
Let $n\in\mathbb{Z}_{>0}$, $A=[a_{jk}],B=[b_{jk}]\in\operatorname{Mat}(n,n,\mathbb{F})$, and $r\in\mathbb{F}$. Then
\begin{enumerate}[label=(\roman*)]
\setlength{\itemsep}{0pt}
\item $\operatorname{tr}(rA)=r\operatorname{tr}(A)$;
\item $\operatorname{tr}(A+B)=\operatorname{tr}(A)+\operatorname{tr}(B)$.
\end{enumerate}
\begin{proof}
\begin{enumerate}[label=(\roman*)]
\setlength{\itemsep}{0pt}
\item $\operatorname{tr}(rA)=\sum\limits_{j=1}^{n}ra_{jj}=r\sum\limits_{j=1}^{n}a_{jj}=r\operatorname{tr}(A)$ by Definition \ref{def:matrix-trace-and-determinant}.
\item $\operatorname{tr}(A+B)=\sum\limits_{j=1}^{n}(a_{jj}+b_{jj})=\sum\limits_{j=1}^{n}a_{jj}+\sum\limits_{j=1}^{n}b_{jj}=\operatorname{tr}(A)+\operatorname{tr}(B)$ by Definition \ref{def:matrix-trace-and-determinant}.
\end{enumerate}
\hlt{Structural Commentary}: Each identity reduces trace to a finite diagonal sum.\\
\hlt{Key Insight}: Finite summation commutes with scalar multiplication and addition.
\end{proof}
\end{proposition}

\begin{example} \label{ex:frobenius-trace-riesz-vector} \hlt{Trace Functional and the Frobenius Inner Product}\\
Let $n\in\mathbb{Z}_{>0}$ and $\mathbb{F}\in\{\mathbb{R},\mathbb{C}\}$, and equip $\operatorname{Mat}(n,n,\mathbb{F})$ with the Frobenius inner product from Example \ref{ex:function-and-matrix-inner-products}.\\
Map $A\mapsto\operatorname{tr}(A)$ is a linear functional by Proposition \ref{prop:trace-arithmetic}. Riesz vector is $I_n$ as $(\forall\,A=[a_{jk}]\in\operatorname{Mat}(n,n,\mathbb{F}))\,\innerprod[F]{A,I_n}=\sum\limits_{j=1}^{n}a_{jj}=\operatorname{tr}(A)$.
\end{example}

\begin{proposition} \label{prop:trace-cyclicity-two-factors} \hlt{Trace Is Cyclic for Two Factors}\\
Let $m,n\in\mathbb{Z}_{>0}$, $A=[a_{jk}]\in\operatorname{Mat}(m,n,\mathbb{F})$, and $B=[b_{kj}]\in\operatorname{Mat}(n,m,\mathbb{F})$. Then $\operatorname{tr}(AB)=\operatorname{tr}(BA)$.
\begin{proof}
$\operatorname{tr}(AB)=\sum\limits_{j=1}^{m}\sum\limits_{k=1}^{n}a_{jk}b_{kj}=\sum\limits_{k=1}^{n}\sum\limits_{j=1}^{m}b_{kj}a_{jk}=\operatorname{tr}(BA)$ by commutativity of $\mathbb{F}$ and interchange of finite sums.\\
\hlt{Structural Commentary}: Entry expansion identifies both traces after interchanging the finite indices.\\
\hlt{Key Insight}: Both traces sum the scalars $a_{jk}b_{kj}$.
\end{proof}
\end{proposition}

\begin{proposition} \label{prop:trace-cyclic-products} \hlt{Cyclic Products}\\
Let $m\in\mathbb{Z}_{\geq2}$ and $n_1,\ldots,n_m\in\mathbb{Z}_{>0}$, and set $n_{m+1}\defeq n_1$. Suppose $(\forall\,j\in\{1,\ldots,m\})\,A_j\in\operatorname{Mat}(n_j,n_{j+1},\mathbb{F})$. Then:
\begin{enumerate}[label=(\roman*)]
\setlength{\itemsep}{0pt}
\item $\operatorname{tr}(A_1\cdots A_m)=\operatorname{tr}(A_mA_1\cdots A_{m-1})$;
\item $(\forall\,n\in\mathbb{Z}_{>0})(\forall\,A,B,C\in\operatorname{Mat}(n,n,\mathbb{F}))\,\operatorname{tr}(ABC)=\operatorname{tr}(BCA)=\operatorname{tr}(CAB)$.
\end{enumerate}
\begin{proof}
\begin{enumerate}[label=(\roman*)]
\setlength{\itemsep}{0pt}
\item Set $X\defeq A_1\cdots A_{m-1}\in\operatorname{Mat}(n_1,n_m,\mathbb{F})$. Then $\operatorname{tr}(A_1\cdots A_m)=\operatorname{tr}(XA_m)=\operatorname{tr}(A_mX)$ by Proposition \ref{prop:trace-cyclicity-two-factors}.
\item Fix $n\in\mathbb{Z}_{>0}$ and $A,B,C\in\operatorname{Mat}(n,n,\mathbb{F})$. Then $\operatorname{tr}(ABC)=\operatorname{tr}(CAB)=\operatorname{tr}(BCA)$ by two applications of part (i).
\end{enumerate}
\hlt{Structural Commentary}: Two-factor cyclicity proves one general rotation, whose repeated use proves the specialization.\\
\hlt{Key Insight}: Grouping $A_1\cdots A_{m-1}$ reduces rotation to two factors.
\end{proof}
\end{proposition}

\begin{example} \label{ex:trace-cyclicity-not-arbitrary-permutation} \hlt{Cyclicity Is Not Arbitrary Commutativity}\\
Let $A=E_{1,2}$, $B=E_{2,1}$, and $C=E_{1,1}$ in $\operatorname{Mat}(2,2,\mathbb{F})$. Then $\operatorname{tr}(ABC)=1$, while $\operatorname{tr}(ACB)=0$.\\
Thus trace permits cyclic rotations of factors, but not arbitrary rearrangements.
\end{example}

For real $2\times2$ matrices, Figure \ref{fig:determinant-column-geometry} illustrates the column geometry behind Theorem \ref{thm:determinant-basic-laws}.

\begin{illustrationfigure}[H]
  \centering
  \begin{tikzpicture}[
  >=Stealth,
  every node/.style={font=\small},
  axis/.style={->, gray!70},
  vec/.style={->, line width=1pt},
  area/.style={fill=blue!12, draw=blue!45}
]
  \draw[axis] (-0.17,0) -- (2.75,0);
  \draw[axis] (0,-0.17) -- (0,1.90);

  \coordinate (O) at (0,0);
  \coordinate (u) at (2.06,0.39);
  \coordinate (v) at (0.69,1.50);
  \coordinate (uv) at ($(u)+(v)$);

  \draw[area] (O) -- (u) -- (uv) -- (v) -- cycle;
  \draw[vec, blue!75!black] (O) -- (u) node[midway, below=3pt] {$\mathbf{a}_1$};
  \draw[vec, red!75!black] (O) -- (v) node[midway, left=2pt] {$\mathbf{a}_2$};
  \draw[gray!55] (u) -- (uv) -- (v);
\end{tikzpicture}

  \caption{The columns $\mathbf{a}_1,\mathbf{a}_2\in\mathbb{R}^2$ span a parallelogram whose signed area is their determinant.}
  \label{fig:determinant-column-geometry}
\end{illustrationfigure}

\begin{computationmethod} \label{meth:determinant-by-structure} \hlt{Determinant by Structure}\\
Given $n\in\mathbb{Z}_{>0}$ and $A\in\operatorname{Mat}(n,n,\mathbb{F})$, the output is $\det A$, computed by the first reduction that lowers the calculation:
\begin{enumerate}[label=(\roman*)]
\setlength{\itemsep}{0pt}
\item if $A$ is triangular, multiply its diagonal entries;
\item if the columns are linearly dependent, including when two coincide, return $0$;
\item if a column decomposition exposes triangular or dependent summands, expand that column by linearity;
\item if a column permutation exposes useful pivots, track the sign of the permutation;
\item if $A=BC$ and the factors have simpler determinants, compute $\det B\det C$.
\end{enumerate}
Each reduction and its output are valid by Theorem \ref{thm:determinant-basic-laws}; verify that every column swap has been counted exactly once.
\end{computationmethod}

For real linear maps, Figure \ref{fig:matrix-determinant-volume-scaling} summarizes the geometric scaling interpretation behind the determinant-by-structure method.

\begin{illustrationfigure}[H]
  \centering
  \begin{tikzpicture}[
  >=Stealth,
  every node/.style={font=\small},
  axis/.style={->, gray!70},
  vec/.style={->, line width=1pt},
  areaA/.style={fill=blue!12, draw=blue!50},
  areaB/.style={fill=red!12, draw=red!55},
  arrow/.style={->, line width=0.85pt, blue!70!black}
]
  \begin{scope}
    \draw[axis] (-0.2,0) -- (2.4,0);
    \draw[axis] (0,-0.2) -- (0,2.0);
    \draw[areaA] (0,0) rectangle (1.35,1.05);
    \node at (0.68,0.52) {$\Omega$};
  \end{scope}

  \draw[arrow] (2.7,0.75) -- (4.1,0.75) node[midway, above=4pt] {$T$};

  \begin{scope}[shift={(4.5,0)}]
    \draw[axis] (-0.2,0) -- (3.0,0);
    \draw[axis] (0,-0.2) -- (0,2.2);
    \coordinate (O) at (0,0);
    \coordinate (a) at (2.25,0.35);
    \coordinate (b) at (0.65,1.55);
    \coordinate (ab) at ($(a)+(b)$);
    \draw[areaB] (O) -- (a) -- (ab) -- (b) -- cycle;
    \draw[vec, blue!75!black] (O) -- (a);
    \draw[vec, red!75!black] (O) -- (b);
    \node at (1.45,0.95) {$T(\Omega)$};
  \end{scope}
\end{tikzpicture}

  \caption{A real linear map scales oriented volume by its determinant; singular maps collapse volume to zero.}
  \label{fig:matrix-determinant-volume-scaling}
\end{illustrationfigure}

\begin{definition} \label{def:schur-complements} \hlt{Schur Complements}\\
Let $r,s\in\mathbb{Z}_{>0}$, $A\in\operatorname{Mat}(r,r,\mathbb{F})$, $B\in\operatorname{Mat}(r,s,\mathbb{F})$, $C\in\operatorname{Mat}(s,r,\mathbb{F})$, and $D\in\operatorname{Mat}(s,s,\mathbb{F})$. Set $M\defeq\left(\begin{smallmatrix}A&B\\C&D\end{smallmatrix}\right)$.
\begin{enumerate}[label=(\roman*)]
\setlength{\itemsep}{0pt}
\item If $A$ is invertible, the \hlt{Schur complement of $A$ in $M$} is $M/A\defeq D-CA^{-1}B$.
\item If $D$ is invertible, the \hlt{Schur complement of $D$ in $M$} is $M/D\defeq A-BD^{-1}C$.
\end{enumerate}
\end{definition}

\begin{proposition} \label{prop:schur-complement-determinant-formula} \hlt{Schur Complement Determinant Formula}\\
Let $r,s\in\mathbb{Z}_{>0}$, $A\in\operatorname{Mat}(r,r,\mathbb{F})$, $B\in\operatorname{Mat}(r,s,\mathbb{F})$, $C\in\operatorname{Mat}(s,r,\mathbb{F})$, $D\in\operatorname{Mat}(s,s,\mathbb{F})$, and $M\defeq\left(\begin{smallmatrix}A&B\\C&D\end{smallmatrix}\right)$.\\
If $A$ is invertible, then $\det M=\det(A)\det(M/A)$.
\begin{proof}
Set $E\defeq\left(\begin{smallmatrix}I_r&0\\-CA^{-1}&I_s\end{smallmatrix}\right)$. Then $EM=\left(\begin{smallmatrix}A&B\\0&M/A\end{smallmatrix}\right)$; $\det E=1$ by Theorem \ref{thm:determinant-basic-laws}(v).\\
Thus $\det M=\det(EM)=\det(A)\det(M/A)$ by Theorems \ref{thm:determinant-basic-laws}(vi) and \ref{thm:block-upper-triangular-determinant}.\\
\hlt{Structural Commentary}: Unit block elimination replaces $C$ by zero and exposes $M/A$.\\
\hlt{Key Insight}: The elimination factor has determinant one.
\end{proof}
\end{proposition}

\begin{proposition} \label{prop:schur-complement-lower-right-determinant} \hlt{Schur Complement of the Lower-Right Block}\\
Let $r,s\in\mathbb{Z}_{>0}$, $A\in\operatorname{Mat}(r,r,\mathbb{F})$, $B\in\operatorname{Mat}(r,s,\mathbb{F})$, $C\in\operatorname{Mat}(s,r,\mathbb{F})$, $D\in\operatorname{Mat}(s,s,\mathbb{F})$, and $M\defeq\left(\begin{smallmatrix}A&B\\C&D\end{smallmatrix}\right)$.\\
If $D$ is invertible, then $\det M=\det(D)\det(M/D)$.
\begin{proof}
Set $E\defeq\left(\begin{smallmatrix}I_r&-BD^{-1}\\0&I_s\end{smallmatrix}\right)$. Then $EM=\left(\begin{smallmatrix}M/D&0\\C&D\end{smallmatrix}\right)$; $\det E=1$ by Theorem \ref{thm:determinant-basic-laws}(v). The transpose of $EM$ is block upper triangular.\\
Thus $\det M=\det(EM)=\det(M/D)\det D$ by Theorems \ref{thm:determinant-basic-laws}(vi),(ix) and \ref{thm:block-upper-triangular-determinant}.\\
\hlt{Structural Commentary}: Unit block elimination replaces $B$ by zero before transposing the triangular form.\\
\hlt{Key Insight}: Transposition converts the lower block-triangular product to the proved upper case.
\end{proof}
\end{proposition}

\begin{proposition} \label{prop:bordered-determinant-formulas} \hlt{Bordered Determinant Formulas}\\
Let $n\in\mathbb{Z}_{>0}$, $A\in\operatorname{Mat}(n,n,\mathbb{F})$, $\mathbf{x},\mathbf{y}\in\mathbb{F}^n$, and $c\in\mathbb{F}$.
\begin{enumerate}[label=(\roman*)]
\setlength{\itemsep}{0pt}
\item If $c\neq0$, then $\det\left(\begin{smallmatrix}c&\mathbf{x}^t\\\mathbf{y}&A\end{smallmatrix}\right)=c\det\left(A-c^{-1}\mathbf{y}\mathbf{x}^t\right)$.
\item If $A$ is invertible, then $\det\left(\begin{smallmatrix}A&\mathbf{x}\\\mathbf{y}^t&c\end{smallmatrix}\right)=\left(c-\mathbf{y}^tA^{-1}\mathbf{x}\right)\det A$.
\end{enumerate}
\begin{proof}
\begin{enumerate}[label=(\roman*)]
\setlength{\itemsep}{0pt}
\item Claim (i) holds by Proposition \ref{prop:schur-complement-determinant-formula} with upper-left block $[c]$ and $M/[c]=A-c^{-1}\mathbf{y}\mathbf{x}^t$.
\item Claim (ii) holds by Proposition \ref{prop:schur-complement-determinant-formula} with upper-left block $A$, $M/A=c-\mathbf{y}^tA^{-1}\mathbf{x}$, and commutativity of $\mathbb{F}$.
\end{enumerate}
\hlt{Structural Commentary}: Each border orientation specializes the upper-left Schur determinant formula.\\
\hlt{Key Insight}: A scalar block turns its Schur complement into a rank-one correction.
\end{proof}
\end{proposition}

\begin{theorem} \label{thm:sylvester-determinant-identity} \hlt{Sylvester Determinant Identity}\\
Let $m,n\in\mathbb{Z}_{>0}$, $A\in\operatorname{Mat}(n,m,\mathbb{F})$, and $B\in\operatorname{Mat}(m,n,\mathbb{F})$. Then $\det(I_n+AB)=\det(I_m+BA)$.
\begin{proof}
Set $M\defeq\left(\begin{smallmatrix}I_n&-A\\B&I_m\end{smallmatrix}\right)$. Then $\det M=\det(I_n+AB)$ by Proposition \ref{prop:schur-complement-lower-right-determinant}, while $\det M=\det(I_m+BA)$ by Proposition \ref{prop:schur-complement-determinant-formula}.\\
\hlt{Structural Commentary}: The two Schur determinant formulas evaluate one block matrix from opposite corners.\\
\hlt{Key Insight}: Both identity blocks have determinant one.
\end{proof}
\end{theorem}

\begin{theorem} \label{thm:ab-ba-nonzero-eigenvalues} \hlt{Nonzero Eigenvalues of $AB$ and $BA$}\\
Let $m,n\in\mathbb{Z}_{>0}$ with $n\geq m$, let $A\in\operatorname{Mat}(m,n,\mathbb{C})$, and let $B\in\operatorname{Mat}(n,m,\mathbb{C})$.
\begin{enumerate}[label=(\roman*)]
\setlength{\itemsep}{0pt}
\item $(\forall z\in\mathbb{C})\,p_{BA}(z)=z^{n-m}p_{AB}(z)$;
\item $(\forall\lambda\in\mathbb{C}\setminus\{0\})\,m_{BA}(\lambda)=m_{AB}(\lambda)$;
\item $m_{BA}(0)=m_{AB}(0)+n-m$.
\end{enumerate}
\begin{proof}
\begin{enumerate}[label=(\roman*)]
\setlength{\itemsep}{0pt}
\item For $z\neq0$, $p_{BA}(z)=z^n\det(I_n-z^{-1}BA)=z^n\det(I_m-z^{-1}AB)=z^{n-m}p_{AB}(z)$ by Theorem \ref{thm:sylvester-determinant-identity} with $-z^{-1}B$ and $A$.\\
If $n=m$, then $p_{BA}(0)=(-1)^n\det(BA)=(-1)^n\det(AB)=p_{AB}(0)$ by Theorem \ref{thm:determinant-basic-laws}(vi) and commutativity of $\mathbb{C}$.\\
If $n>m$, then $\operatorname{col}(BA)\subseteq\operatorname{col}(B)$ because $(\forall\mathbf{x}\in\mathbb{C}^n)\,BA\mathbf{x}=B(A\mathbf{x})$.\\
Thus $\operatorname{rank}(BA)\leq\operatorname{rank}(B)\leq m<n$ by Definition \ref{def:column-space-matrix-rank-nullity}, Proposition \ref{prop:dimension-bounds-and-basis-criteria}(i), and Corollary \ref{cor:matrix-rank-nullity}(i).\\
$BA$ singular by Corollary \ref{cor:square-matrix-invertible-iff-full-rank}. Then $p_{BA}(0)=(-1)^n\det(BA)=0=0^{n-m}p_{AB}(0)$ by Definition \ref{def:matrix-characteristic-polynomial}, Theorem \ref{thm:determinant-basic-laws}(vii).
\item $m_{BA}(\lambda)=m_{AB}(\lambda)$ for $\lambda\neq0$ by part (i), because $z^{n-m}$ has multiplicity zero at $\lambda$.
\item $m_{BA}(0)=m_{AB}(0)+n-m$ by part (i), because $z^{n-m}$ has multiplicity $n-m$ at $0$.
\end{enumerate}
\hlt{Structural Commentary}: Sylvester transfer proves the polynomial identity before root orders split into zero and nonzero cases.\\
\hlt{Key Insight}: The size difference appears only as the factor $z^{n-m}$.
\end{proof}
\end{theorem}

\begin{theorem} \label{thm:similarity-preserves-trace-and-determinant} \hlt{Trace and Determinant Are Similarity Invariants}\\
Let $n\in\mathbb{Z}_{>0}$ and let $A,B\in\operatorname{Mat}(n,n,\mathbb{F})$ be similar.
\begin{enumerate}[label=(\roman*)]
\setlength{\itemsep}{0pt}
\item $\operatorname{tr}A=\operatorname{tr}B$;
\item $\det A=\det B$.
\end{enumerate}
\begin{proof}
Choose an invertible $P$ with $B=PAP^{-1}$ by Definition \ref{def:similar-matrices}.
\begin{enumerate}[label=(\roman*)]
\setlength{\itemsep}{0pt}
\item $\operatorname{tr}B=\operatorname{tr}\bigl(P(AP^{-1})\bigr)=\operatorname{tr}\bigl((AP^{-1})P\bigr)=\operatorname{tr}A$ by Proposition \ref{prop:trace-cyclicity-two-factors}.
\item $\det B=\det(P)\det(A)\det(P^{-1})=\det A$ by Theorem \ref{thm:determinant-basic-laws}(vi),(viii).
\end{enumerate}
\hlt{Structural Commentary}: Cyclicity and multiplicativity remove the conjugating factors from the two invariants.\\
\hlt{Key Insight}: In both calculations, $P^{-1}P=I_n$.
\end{proof}
\end{theorem}

\begin{proposition} \label{prop:trace-determinant-from-eigenvalues} \hlt{Trace and Determinant from Eigenvalues}\\
Let $n\in\mathbb{Z}_{>0}$, let $A\in\operatorname{Mat}(n,n,\mathbb{C})$, and let $\lambda_1,\ldots,\lambda_n$ list its eigenvalues with algebraic multiplicity.
\begin{enumerate}[label=(\roman*)]
\setlength{\itemsep}{0pt}
\item $\operatorname{tr}A=\sum\limits_{j=1}^{n}\lambda_j$;
\item $\det A=\prod\limits_{j=1}^{n}\lambda_j$.
\end{enumerate}
\begin{proof}
$p_A(z)=\prod\limits_{j=1}^{n}(z-\lambda_j)$ by the algebraic-multiplicity list and monicity.\\
The identity $p_A(z)=z^n-(\operatorname{tr}A)z^{n-1}+\cdots+(-1)^n\det A$ holds by Theorem \ref{thm:matrix-characteristic-polynomial-basic-coefficients}.
\begin{enumerate}[label=(\roman*)]
\setlength{\itemsep}{0pt}
\item $\operatorname{tr}A=\sum\limits_{j=1}^{n}\lambda_j$ by comparison of the $z^{n-1}$ coefficients.
\item $\det A=\prod\limits_{j=1}^{n}\lambda_j$ by comparison of the constant coefficients.
\end{enumerate}
\hlt{Structural Commentary}: Two coefficient comparisons transfer the root factorization to trace and determinant.\\
\hlt{Key Insight}: The next-to-leading and constant coefficients encode the root sum and product.
\end{proof}
\end{proposition}

\begin{proposition} \label{prop:transpose-conjugate-adjoint-characteristic-polynomial} \hlt{Transpose, Conjugate, and Adjoint Spectra}\\
Let $n\in\mathbb{Z}_{>0}$, $A\in\operatorname{Mat}(n,n,\mathbb{C})$, and $\lambda\in\operatorname{spec}(A)$.
\begin{enumerate}[label=(\roman*)]
\setlength{\itemsep}{0pt}
\item $p_{A^t}=p_A$;
\item $(\forall z\in\mathbb{C})\,p_{\overline{A}}(z)=\overline{p_A(\overline{z})}$;
\item $(\forall z\in\mathbb{C})\,p_{A^*}(z)=\overline{p_A(\overline{z})}$;
\item $m_{A^t}(\lambda)=m_A(\lambda)$;
\item $m_{\overline{A}}(\overline{\lambda})=m_A(\lambda)$;
\item $m_{A^*}(\overline{\lambda})=m_A(\lambda)$.
\end{enumerate}
\begin{proof}
\begin{enumerate}[label=(\roman*)]
\setlength{\itemsep}{0pt}
\item $(\forall z\in\mathbb{C})\,p_{A^t}(z)=\det((zI_n-A)^t)=p_A(z)$ by Theorem \ref{thm:determinant-basic-laws}(ix).
\item $(\forall X\in\operatorname{Mat}(n,n,\mathbb{C}))\,\det(\overline{X})=\overline{\det X}$ by conjugation of each term in the permutation formula.\\
Hence $(\forall z\in\mathbb{C})\,p_{\overline{A}}(z)=\det(\overline{\overline{z}I_n-A})=\overline{p_A(\overline{z})}$.
\item $(\forall z\in\mathbb{C})\,p_{A^*}(z)=p_{\overline{A}}(z)=\overline{p_A(\overline{z})}$ by $A^*=\overline{A}^{\,t}$ and parts (i)--(ii).
\item $m_{A^t}(\lambda)=m_A(\lambda)$ by part (i).
\item Set $k\defeq m_A(\lambda)$. Choose $q\in\mathbb{C}[z]$ satisfying $p_A(z)=(z-\lambda)^kq(z)$ and $q(\lambda)\neq0$ by Definition \ref{def:matrix-algebraic-multiplicity}.\\
Identity $p_{\overline{A}}(z)=(z-\overline{\lambda})^k\overline{q(\overline{z})}$ holds by part (ii), while $\overline{q(\lambda)}\neq0$ by choice of $q$. Thus $m_{\overline{A}}(\overline{\lambda})=k=m_A(\lambda)$ by Definition \ref{def:matrix-algebraic-multiplicity}.
\item $m_{A^*}(\overline{\lambda})=m_A(\lambda)$ by parts (iii) and (v).
\end{enumerate}
\hlt{Structural Commentary}: Determinant transposition preserves factors, while entrywise conjugation transports each root.\\
\hlt{Key Insight}: The adjoint adds no spectral change beyond complex conjugation.
\end{proof}
\end{proposition}

\begin{perspective} \label{per:similarity-invariants-and-canonical-forms} \hlt{Similarity Invariants and Canonical Forms}\\
For square complex matrices, trace and determinant are similarity invariants by Theorem \ref{thm:similarity-preserves-trace-and-determinant}.\\
The characteristic and minimal polynomials are similarity invariants by Theorems \ref{thm:similarity-preserves-characteristic-polynomial} and \ref{thm:similarity-preserves-minimal-polynomial}.\\
Set $J\defeq J_3(0)\oplus J_3(0)$ and $K\defeq J_3(0)\oplus J_2(0)\oplus J_1(0)$.\\
Both have characteristic polynomial $z^6$ by Proposition \ref{prop:upper-triangular-characteristic-polynomial}(i).\\
Both are nilpotent of index $3$ by Lemma \ref{lem:nilpotent-jordan-block-power-rank} and Proposition \ref{prop:direct-sum-matrix-algebra}(iii).\\
Hence $\mu_J=\mu_K=z^3$ by Theorem \ref{thm:operator-annihilator-iff-multiple-minimal-polynomial} and Definition \ref{def:minimal-polynomial-matrix}.\\
They are not similar by Corollary \ref{cor:jordan-form-unique-up-to-block-permutation}(i), because their Jordan block multisets differ.\\
Thus Theorem \ref{thm:complex-jordan-form} and that corollary make the Jordan block multiset a complete similarity invariant.
\end{perspective}

\begin{theorem} \label{thm:operator-matrix-trace-determinant-basis-independent} \hlt{Trace and Determinant of an Operator Are Basis-Independent}\\
Let $\mathcal{V}\neq\{\mathbf{0}\}$ be finite-dimensional, let $T\in\LL(\mathcal{V})$, and let $\beta,\gamma$ be bases of $\mathcal{V}$.
\begin{enumerate}[label=(\roman*)]
\setlength{\itemsep}{0pt}
\item $\operatorname{tr}\bigl([T]_{\beta\leftarrow\beta}\bigr)=\operatorname{tr}\bigl([T]_{\gamma\leftarrow\gamma}\bigr)$;
\item $\det\bigl([T]_{\beta\leftarrow\beta}\bigr)=\det\bigl([T]_{\gamma\leftarrow\gamma}\bigr)$.
\end{enumerate}
\begin{proof}
$[T]_{\gamma\leftarrow\gamma}=C_{\gamma\leftarrow\beta}[T]_{\beta\leftarrow\beta}C_{\beta\leftarrow\gamma}$ by Theorem \ref{thm:change-of-basis-formula}(i). Thus the representing matrices are similar by Definition \ref{def:similar-matrices}.
\begin{enumerate}[label=(\roman*)]
\setlength{\itemsep}{0pt}
\item The traces are equal by Theorem \ref{thm:similarity-preserves-trace-and-determinant}(i).
\item The determinants are equal by Theorem \ref{thm:similarity-preserves-trace-and-determinant}(ii).
\end{enumerate}
\hlt{Structural Commentary}: Change of basis reduces both claims to the corresponding similarity invariants.\\
\hlt{Key Insight}: Every representing matrix lies in one conjugacy class.
\end{proof}
\end{theorem}

Figure \ref{fig:operator-trace-determinant-basis-invariance} summarizes the conjugation relation that makes operator trace and determinant basis-independent.

\begin{illustrationfigure}[H]
  \centering
  \begin{tikzpicture}[
  >=Stealth,
  every node/.style={font=\small},
  box/.style={draw=gray!55, rounded corners=3pt, minimum width=2.5cm, minimum height=0.84cm, align=center, fill=gray!4},
  map/.style={->, line width=0.8pt, blue!65!black}
]
  \node[box] (beta) at (0,0) {$[T]_{\beta\leftarrow\beta}=A$};
  \node[box] (gamma) at (4.2,0) {$[T]_{\gamma\leftarrow\gamma}=PAP^{-1}$};

  \draw[map] (beta) -- (gamma)
    node[midway, above] {$P=C_{\gamma\leftarrow\beta}$};
\end{tikzpicture}

  \caption{Basis changes conjugate matrices; trace and determinant are unchanged.}
  \label{fig:operator-trace-determinant-basis-invariance}
\end{illustrationfigure}

\begin{definition} \label{def:operator-trace-and-determinant} \hlt{Trace and Determinant of an Operator}\\
Let $\mathcal{V}\neq\{\mathbf{0}\}$ be finite-dimensional, let $T\in\LL(\mathcal{V})$, and choose a basis $\beta$ of $\mathcal{V}$.
\begin{enumerate}[label=(\roman*)]
\setlength{\itemsep}{0pt}
\item The \hlt{trace} of $T$ is $\operatorname{tr}T\defeq\operatorname{tr}\bigl([T]_{\beta\leftarrow\beta}\bigr)$.
\item The \hlt{determinant} of $T$ is $\det T\defeq\det\bigl([T]_{\beta\leftarrow\beta}\bigr)$.
\end{enumerate}
Both quantities are independent of $\beta$ by Theorem \ref{thm:operator-matrix-trace-determinant-basis-independent}.
\end{definition}

\begin{proposition} \label{prop:isometry-determinant-absolute-one} \hlt{Unitary Matrices and Isometries Have Determinant of Modulus One}\\
Let $\mathbb{F}\in\{\mathbb{R},\mathbb{C}\}$.
\begin{enumerate}[label=(\roman*)]
\setlength{\itemsep}{0pt}
\item If $n\in\mathbb{Z}_{>0}$ and $U\in\operatorname{Mat}(n,n,\mathbb{F})$ is unitary, then $\abs{\det U}=1$.
\item If $\mathcal{V}\neq\{\mathbf{0}\}$ is a finite-dimensional inner-product space over $\mathbb{F}$ and $S\in\LL(\mathcal{V})$ is an isometry, then $\abs{\det S}=1$.
\end{enumerate}
\begin{proof}
\begin{enumerate}[label=(\roman*)]
\setlength{\itemsep}{0pt}
\item $\det(U^*)=\overline{\det U}$ by the permutation formula and Theorem \ref{thm:determinant-basic-laws}(ix).\\
Taking determinants in $U^*U=I_n$, $1=\overline{\det U}\det U=\abs{\det U}^2$ by Theorem \ref{thm:determinant-basic-laws}(i),(vi), so $\abs{\det U}=1$.
\item Set $n\defeq\dim\mathcal{V}$. Choose an orthonormal basis $\beta$ by Corollary \ref{cor:finite-dimensional-inner-product-space-has-orthonormal-basis} and set $U\defeq[S]_{\beta\leftarrow\beta}$.\\
Then $U^*U=I_n$ by Theorems \ref{thm:finite-dimensional-isometry-iff-adjoint-product-identity} and \ref{thm:matrix-of-adjoint-is-conjugate-transpose} and Proposition \ref{prop:composition-and-matrix-multiplication}. Thus $\abs{\det S}=\abs{\det U}=1$ by part (i) and Definition \ref{def:operator-trace-and-determinant}.
\end{enumerate}
\hlt{Structural Commentary}: The matrix identity is proved first, then orthonormal coordinates transfer it to operators.\\
\hlt{Key Insight}: Determinants convert $U^*U=I_n$ into $\abs{\det U}^2=1$.
\end{proof}
\end{proposition}

\begin{proposition} \label{prop:operator-trace-determinant-characteristic-coefficients} \hlt{Trace and Determinant from $q_T$}\\
Let $\mathcal{V}\neq\{\mathbf{0}\}$ be finite-dimensional over $\mathbb{F}\in\{\mathbb{R},\mathbb{C}\}$, let $n\defeq\dim\mathcal{V}$, let $T\in\LL(\mathcal{V})$, and write $q_T(z)=z^n+c_{n-1}z^{n-1}+\cdots+c_1z+c_0$.
\begin{enumerate}[label=(\roman*)]
\setlength{\itemsep}{0pt}
\item $\operatorname{tr}T=-c_{n-1}$;
\item $\det T=(-1)^nc_0$;
\item if $\mathbb{F}=\mathbb{C}$ and $\lambda_1,\ldots,\lambda_n$ list the eigenvalues of $T$ with algebraic multiplicity, then $\operatorname{tr}T=\sum\limits_{j=1}^{n}\lambda_j$;
\item under the hypotheses of (iii), $\det T=\prod\limits_{j=1}^{n}\lambda_j$.
\end{enumerate}
\begin{proof}
Choose a basis $\beta$ and set $A\defeq[T]_{\beta\leftarrow\beta}$.\\
Then $q_T=p_A$ by Definition \ref{def:operator-characteristic-polynomial}, while $\operatorname{tr}T=\operatorname{tr}A$ and $\det T=\det A$ by Definition \ref{def:operator-trace-and-determinant}.
\begin{enumerate}[label=(\roman*)]
\setlength{\itemsep}{0pt}
\item $\operatorname{tr}T=-c_{n-1}$ by Theorem \ref{thm:matrix-characteristic-polynomial-basic-coefficients}.
\item $\det T=(-1)^nc_0$ by Theorem \ref{thm:matrix-characteristic-polynomial-basic-coefficients}.
\item Under the hypotheses of (iii), $\lambda_1,\ldots,\lambda_n$ also list the eigenvalues of $A$ with algebraic multiplicity by $q_T=p_A$.\\
Hence $\operatorname{tr}T=\sum\limits_{j=1}^{n}\lambda_j$ by Proposition \ref{prop:trace-determinant-from-eigenvalues}(i).
\item Under the same hypotheses, $\det T=\prod\limits_{j=1}^{n}\lambda_j$ by Proposition \ref{prop:trace-determinant-from-eigenvalues}(ii) with the list from part (iii).
\end{enumerate}
\hlt{Structural Commentary}: A basis identifies $q_T$ with $p_A$, after which matrix coefficient formulas transfer claimwise.\\
\hlt{Key Insight}: Operator trace and determinant are the corresponding matrix invariants in any basis.
\end{proof}
\end{proposition}
